\begin{table}[t]
\centering
\caption{Ejemplo num\'erico para comprobar una implementaci\'on de FLC-O: pertenencias de $G=2$, $R=5$ y $V=10$\,s con los nudos de la Tabla~\ref{tab:sparametros}; la salida resulta $u=-1.4478$ (seguir en verde).}
\label{tab:sejemplo}
\scriptsize
\begin{tabular}{@{}lrrr@{}}
\toprule
Entrada & bajo & medio & alto \\
\midrule
G = 2 & 0.0000 & 0.7512 & 0.2488 \\
R = 5 & 0.3903 & 0.6097 & 0.0000 \\
V = 10 & 0.7965 & 0.2035 & 0.0000 \\
\bottomrule
\end{tabular}
\fuente{c\'alculo con las Ecs. del art\'iculo (opt\_flc2.py).}
\end{table}
